\thispagestyle{empty}
\chapter*{译注本说明}
\addcontentsline{toc}{chapter}{译注本说明}
\markboth{译注本说明}{译注本说明}

本书是 D.~O.~Orlov 综述《凝聚层的导出范畴及其之间的等价》（\emph{Derived categories of
coherent sheaves and equivalences between them}, Russian Math. Surveys \textbf{58}:3 (2003),
511--591）的中文译注本。原文是 2003 年作者为《俄罗斯数学综览》所写的综述，系统总结了
Bondal--Orlov、Orlov 本人以及 Bridgeland、Mukai 等人关于凝聚层导出范畴的工作，
其主定理（每个导出范畴之间的等价都由一个核对象给出）是这一方向的奠基性结果之一。

\subsection*{一、翻译底本与页码}

译文依据英文版排版，逐句对照俄文原刊的分页（两者页码相差 422）。为便于与原文对照，
正文源码在每一页的起首记录原书页码（形如 \texttt{\% [原文 p.545]}）；这些标记不进入成品
版面，读者如需按页核对原文，可依附录 D 的定位表查相应位置。原书共 31 个编号公式、
若干交换图；公式编号一律照抄原书号，未重新编号。

\subsection*{二、体例}

\begin{enumerate}[label=（\arabic*）,leftmargin=2.6em,itemsep=2pt,topsep=4pt]
\item \textbf{陈述编号照抄原书}。原书按节编号（如「定理 2.1.8」指第 2 章第 1 节第 8 条），
      本书沿用同一编号体系，不使用 LaTeX 自动编号，以便与原文逐条对应。
\item \textbf{两种批注}。随文可见的「〔译注〕」用于解释概念、交代背景、指出与 nLab 或
      经典文献的对应，每条都以【来源：……】或〔译者按（上下文推断，无外部依据）：……〕收尾，
      来源按 A（原著与经典教材）、B（nLab、Stacks Project，凡单独引用者标「待核」）、
      I（本库既往译本）三级标注。不进入正文的校勘性批注另编为附录 B，专记原文笔误、
      原书排印出入与本译本的取舍。
\item \textbf{原书笔误照录}。凡原文排印有误处，译文保留原样，另在附录 B 中说明，
      不静默改字，以免读者对照原文时对不上。
\item \textbf{术语}。有稳定通行译名者用中文，无通行译名或两种译名并存者保留英文原词
      （如 flop、Fourier--Mukai 变换、Postnikov 系统）。全部术语及其中英文对照见附录 A。
\item \textbf{参考文献}。文献条目照排英文原文，不译，仅将 PDF 文字层中被破坏的重音符号
      与断字恢复为正常排印。
\end{enumerate}

\subsection*{三、附录}

\begin{itemize}[leftmargin=2.6em,itemsep=1pt,topsep=2pt]
\item \textbf{附录 A}\quad 术语对照表（英--中，含依据等级）。
\item \textbf{附录 B}\quad 译注汇编（校勘性批注，按原书页码排列）。
\item \textbf{附录 C}\quad 概念批注：与 nLab 的对照（对书中十余个核心概念给出
      nLab 对应条目、准确定义与在本文论证中的作用，并说明 nLab 视角与
      Orlov 原文用语的差异）。
\item \textbf{附录 D}\quad 定理与定义索引（按原书编号排列，附中文章节定位）。
\end{itemize}

\subsection*{四、致读者}

Orlov 这篇综述的技术密度集中在第 3 章：前半（1--2 章）是标准预备知识与半正交分解的
工具箱，第 3 章把「全忠实 + 有伴随」逐步翻译成「由核对象表示」，第 4--5 章则是这一
主定理在 K3 曲面与阿贝尔簇上的两个精细应用。若读者只关心主定理，可先读引言、
定理 2.1.5、定理 3.2.1 与第 3.4 节；若关心应用，第 4 章的 Mukai 格与第 5 章的
等距同构判据是两条彼此平行、方法相通的主线。每章末的「本章回顾与评注」与每节首的
「本节重点」可供快速定位。

\vspace{1em}
\noindent\hfill 译者\quad 二〇二六年十一月
\clearpage
